\section{Las personas: el lenguaje es una herramienta inventada para generalizar}

Esta sección examina primero la parte de «las personas». Una afirmación central de Yao Shunyu (姚顺雨) es que el lenguaje es una herramienta que los seres humanos inventaron para lograr la generalización, y que esto es más fundamental que cualquier otra cosa. Esta idea explica por qué él partió del lenguaje para trabajar en Agent: el lenguaje no es solo un medio de comunicación, sino también la herramienta con la que los seres humanos comprimen la experiencia, transmiten reglas, describen objetivos, organizan planes y transfieren lo aprendido de un escenario a otro.

\begin{figure}[H]
\centering
\includegraphics[width=0.94\textwidth]{language-generalization.png}
\caption{El lenguaje como herramienta de generalización: los seres humanos comprimen la experiencia con el lenguaje, y el Agent usa el lenguaje para conectar las tareas con el mundo. Diagrama conceptual de elaboración propia, basado en la conversación de 00:02:03--00:07:50.}
\end{figure}

\begin{knowledgebox}{Lectura de la figura: por qué el lenguaje es adecuado como capa intermedia del Agent}
El lenguaje convierte las tareas, las reglas, las restricciones, las descripciones de herramientas y la retroalimentación en símbolos que pueden combinarse. Mediante el lenguaje, el modelo puede comprender los objetivos humanos y también invocar código, API, documentos y otros Agent. Por ello, el lenguaje no es todo el Agent, pero sí es una interfaz general muy potente.
\end{knowledgebox}

\subsection{Una postura fuera del consenso: apostar desde el principio por los Agent de lenguaje}

La sección anterior explicó por qué el lenguaje es una herramienta de generalización; esta examina cómo ese juicio se convirtió en una apuesta de investigación. Yao Shunyu dice que siempre sostuvo una postura fuera del consenso: quería trabajar en Agent y, como primer paso, construirlos sobre modelos de lenguaje. En ese momento, la corriente dominante no necesariamente consideraba que los modelos de lenguaje fueran el mejor camino hacia los Agent; pero después la serie GPT mostró que los modelos de lenguaje tienen un gran potencial de generalización y de servir como interfaz con herramientas, y esa postura fuera del consenso pasó a ser la línea principal de la frontera de investigación.

\begin{importantbox}{El sentido de una Different Bet}
Si solo se copia la línea principal del líder ya establecido, es muy difícil superarlo. Las nuevas oportunidades suelen surgir de una different bet: tareas distintas, entornos distintos, interface distinta, formas de producto distintas. En sus inicios, los Agent de lenguaje fueron justamente una apuesta de este tipo.
\end{importantbox}

\begin{knowledgebox}{Por qué los primeros Agent de lenguaje estaban fuera del consenso}
Los primeros modelos de lenguaje parecían más sistemas de texto que agentes inteligentes capaces de actuar. Para verlos como Agent había que creer tres cosas a la vez: que el lenguaje puede expresar tareas, que el modelo puede generalizar de una tarea a otra y que las herramientas externas pueden convertir los planes expresados en lenguaje en acciones. En ese momento, ninguna de las tres era un consenso evidente.
\end{knowledgebox}

\begin{knowledgebox}{Lectura didáctica de la trayectoria de investigación de Yao Shunyu}
\begin{tabularx}{\textwidth}{p{0.20\textwidth}p{0.34\textwidth}X}
\toprule
Etapa & Problema de interés & Relación con el Agent \\
\midrule
Inicio del doctorado & Usar modelos de lenguaje para juegos y tareas sencillas & Demostrar que el lenguaje puede contener decisiones y estados. \\
Entornos de tareas & Buscar benchmark más realistas y valiosos & Lograr que el Agent no se limite a resolver ejercicios pequeños. \\
Herramientas/código & InterCode, entornos de código, affordance del mundo digital & Darle al Agent manos capaces de ejecutar. \\
Segunda mitad & reward, multi-agent, interface, evolución de sistemas & Pasar de la capacidad de un solo modelo a la capacidad del sistema. \\
\bottomrule
\end{tabularx}
\end{knowledgebox}

\subsection{Resumen de la sección}

La sección sobre «las personas» muestra que la investigación sobre Agent no surgió de la nada a partir de las capacidades de los modelos, sino de una comprensión del lenguaje, de la generalización humana y de la expresión de tareas. El lenguaje es importante porque conecta los objetivos humanos con la acción de las máquinas.
